\paragraph{Real target contractions of other signatures (algebraic at every \(N\)).}
Keep Definition~(A) and the notation of the three preceding paragraphs, but replace \(\eta\) in \(S_{\mathrm b}\) by a real symmetric invertible \(h_0\) with \(10-t\) positive and \(t\) negative eigenvalues, \(S_{\mathrm b}=-\tfrac1{4g_{\mathrm{YM}}^2}\sum_{\mu,\nu,\rho,\sigma}(h_0)_{\mu\rho}(h_0)_{\nu\sigma}\operatorname{Tr}[A_\mu,A_\nu][A_\rho,A_\sigma]\), and let the factor \(w(A)\) multiplying \(\mathrm e^{\mathrm iS_{\mathrm b}}\) be a homogeneous polynomial that does not vanish identically, such as \(1\) or a Pfaffian. On the contour \(A_\mu=\sum_\nu\Lambda_{\mu\nu}X_\nu\) put \(h=\Lambda^{\mathrm T}h_0\Lambda\). The formula for \(S_{\mathrm b}\) at \(N=2\) and the set \(\mathfrak C=\mathfrak C_0\cup(-\mathfrak C_0)\) are then those of the paragraph on linear contours that mix the indices, which is the case \(h_0=\eta\). Since \(S_{\mathrm b}\) is even in \(h_0\), the values \(t\) and \(10-t\) give the same integrand, with \(\mathfrak C_0\) and \(-\mathfrak C_0\) exchanged. At \(N=2\) the condition \(\operatorname{Re}\mathrm iS_{\mathrm b}\le0\) on the contour is \(\operatorname{Im}\det(h|_E)\ge0\) for every real 2-plane \(E\), the determinant being taken in a real basis of \(E\), and for \(N\ge3\) it implies this, by restriction to an \(\mathfrak{su}(2)\) block.

(i) \emph{Real boundary points.} By (ii) of the paragraph on linear contours that mix the indices, \(\operatorname{Re}h\) has at least nine non-negative eigenvalues on the closure of \(\mathfrak C_0\) and at least nine non-positive ones on that of \(-\mathfrak C_0\). So the real invertible forms in the closure of \(\mathfrak C\) are those with \(t\in\{0,1,9,10\}\), and all of these occur: \(\mathfrak C_0\) is invariant under real congruences, \(T^{\mathrm T}T\) is the limit of \(\mathrm e^{\mathrm i\theta}T^{\mathrm T}T\in\mathfrak C_0\) as \(\theta\to0^+\), and \(\eta\) is the limit of the \(\mathcal R_u\). This is an analogue, for the quartic weight of the matrix integral, of the allowable complex metrics proposed in~\cite{KontsevichSegal2021} as the domain to which quantum field theories continue analytically, which contains the Riemannian metrics and has the Lorentzian ones on its boundary. In a real basis diagonalising such a metric the condition is \(\sum_i|\operatorname{Arg}\lambda_i|<\pi\), Eq.~(3.5) of~\cite{Witten2021ComplexMetrics}, which a real metric with \(t\ge2\) violates also in its closed form. Two differences matter. \(\mathfrak C\) is invariant under \(h\mapsto-h\), whence the two classes and the values \(t=9,10\), whereas the allowable domain contains no negative definite metric. And~\cite{Witten2021ComplexMetrics} reaches a Lorentzian metric from the allowable domain in two ways, with time component \(-(1\mp\mathrm i\epsilon)\), which give the Feynman integral and its complex conjugate. For the matrix weight the corresponding pair is \(h\) and \(\bar h\). Since \(\operatorname{Im}\det(\bar h|_E)=-\operatorname{Im}\det(h|_E)\), for \(h\in\mathfrak C\) the form \(\bar h\) damps \(\mathrm e^{-\mathrm iS_{\mathrm b}}\), the opposite sign of the quartic term noted under Definition~(A), whose integral is the complex conjugate. The classes \(\mathfrak C_0\) and \(-\mathfrak C_0\) are not this pair. They are exchanged by the overall sign of the contraction, and \(\eta\) lies in the closure of \(\mathfrak C_0\) only, \(-\eta\) in that of \(-\mathfrak C_0\) only.

(ii) \emph{Rigidity.} Let \(h\) be complex symmetric with \(\operatorname{Im}\det(h|_E)\ge0\) for every real 2-plane \(E\), and let \(R=\operatorname{Re}h\) be invertible with at least two positive and at least two negative eigenvalues. Then \(h\) is real. For real symmetric \(2\times2\) matrices \(r,s\) one has \(\operatorname{Im}\det(r+\mathrm is)=\operatorname{tr}(\operatorname{adj}(r)\,s)\). Write \(S=\operatorname{Im}h\) and split \(\mathbb R^{10}=V_+\oplus V_-\), orthogonally for \(R\), with \(R\) positive definite on \(V_+\) and negative definite on \(V_-\). Let \(\lambda_1\le\dots\le\lambda_p\) be the eigenvalues of \(S\) on \(V_+\) relative to \(R\), and \(\nu_1\le\dots\le\nu_q\) those on \(V_-\) relative to \(-R\). A plane in \(V_+\), in an \(R\)-orthonormal basis, gives \(\lambda_1+\lambda_2\ge0\). A plane in \(V_-\) gives \(\nu_{q-1}+\nu_q\le0\). A plane spanned by \(e\in V_+\) and \(f\in V_-\) with \(R(e,e)=1=-R(f,f)\) gives \(S(f,f)\ge S(e,e)\), so \(\lambda_p\le\nu_1\). Hence \(0\le\lambda_2\le\lambda_p\le\nu_1\le\nu_{q-1}\le0\), and with the two sums all \(\lambda_i\) and \(\nu_j\) vanish. Finally, for \(e_1\in V_+\) with \(R(e_1,e_1)=1\) and \(R\)-orthogonal \(f_1,f_2\in V_-\) with \(R(f_k,f_k)=-1\), the plane spanned by \(\cosh a\,e_1+\sinh a\,f_1\) and \(f_2\) gives \(-\sinh(2a)\,S(e_1,f_1)\ge0\) for all real \(a\), so \(S(e_1,f_1)=0\) and \(S=0\). By continuity of eigenvalues, for \(2\le t\le8\) the real forms with \(t\) negative eigenvalues are therefore open and closed in the set of invertible complex symmetric \(h\) with \(\operatorname{Im}\det(h|_E)\ge0\) for all \(E\). In particular they are not in the closure of \(\mathfrak C\), which gives (i) again for these \(t\).

(iii) \emph{Contours.} Let \(2\le t\le8\). (a) Let \(\Lambda(\tau)\), \(0\le\tau\le1\), be continuous and invertible, with \(\Lambda(0)\) real and \(\operatorname{Re}\mathrm iS_{\mathrm b}\le0\) on every contour of the path. Then \(h(\tau)\) is real for all \(\tau\): it is real and of constant signature up to the supremum of the \(\tau\) at which it is real, and by (ii) beyond it. This holds at every \(N\), for every regulator and with any mass term, which do not enter the condition. (b) Suppose that in a real basis in which the regulator form \(P\) is the identity, \(h_0\) is diagonal with entries \(\pm c\), \(c>0\), as for Definition~(A) with \(h_0=\mathrm{diag}(\pm1)\). With \(M=\Lambda^{\mathrm T}P\Lambda\) and \(\ell_j\) the rows of \(\Lambda\) for the \(t\) negative entries, \(h=c(M-2\sum_j\ell_j\ell_j^{\mathrm T})\), and \(hv=-cMv\) on the span of the \(\Lambda^{-1}e_j\). If \(\operatorname{Re}M\) is positive definite, then \(\operatorname{Re}h\) is positive definite on the real \(x\) with \(\operatorname{Re}\ell_j^{\mathrm T}x=0\) for all \(j\), a subspace of dimension at least \(10-t\), and \(\bar v^{\mathrm T}(\operatorname{Re}h)v=-c\,\bar v^{\mathrm T}(\operatorname{Re}M)v<0\) on that span, so \(\operatorname{Re}h\) has the inertia of \(h_0\), as in (i) of the paragraph on the regulator for \(t=1\). By (ii), every linear contour with \(\operatorname{Re}M\) positive definite and \(\operatorname{Re}\mathrm iS_{\mathrm b}\le0\) then has real \(h\), whether or not a path leads to it. (c) On the contours of (a) and (b), \(S_{\mathrm b}\) is real at every \(N\), so none of them is damped, and at \(\epsilon=0\) the integrand has modulus \(|w(\Lambda X)|\), which is not integrable, since \(w\circ\Lambda\) is a homogeneous polynomial that does not vanish identically. The dominated convergence used for \(t=1\) in the paragraph on the regulated Lorentzian integral is therefore unavailable on every such contour, and on every invertible limit of such contours as \(\epsilon\to0\), whose \(h\) is again real, since a dominating function would also dominate the integrand at \(\epsilon=0\) on the limit. For \(t\in\{0,1\}\) paths ending on damped contours exist: the \(\mathcal R_u\) for \(h_0=\eta\), and \(A=\mathrm e^{\pi\mathrm i\tau/8}X\) for \(h_0=\mathbb 1\), with \(\operatorname{Re}M=\cos(\pi\tau/4)\mathbb 1\) and \(\operatorname{Re}\mathrm iS_{\mathrm b}=-\sin(\pi\tau/2)S_{\mathrm E}(X)\). For \(t\in\{9,10\}\) the same contours serve, by evenness.

(iv) \emph{Fermion content.} A Majorana--Weyl spinor exists in signature \((s,t)\) exactly when \(s-t\equiv0\bmod8\): the Majorana condition and the compatibility of the reality condition with the chiral projection, Eqs.~(3.25) and~(3.28) of~\cite{VanProeyen1999}, together leave these cases, as in Table~2 there. For \(s+t=10\) this means \(t\in\{1,5,9\}\). The content of~\cite{IKKT1997} is one Majorana--Weyl spinor.

So, for linear contours that do not mix the colour components, Definition~(A) and \(\epsilon\to0\) taken at fixed \(N\), two premises fix the signature of the contraction. (P1) The limit is obtained, as for \(t=1\) in the paragraph on the regulated Lorentzian integral, by dominated convergence on a damped linear contour reached from the Hermitian contour by a path as in (iii)(a). (P2) The supplied fermion is a single Majorana--Weyl spinor for \(h_0\). Then \(t\in\{0,1,9,10\}\cap\{1,5,9\}=\{1,9\}\), so \(h_0=\pm T^{\mathrm T}\eta T\) with \(T\) real and invertible. \(S_{\mathrm b}\) does not see the sign, and \(T\) turns the regulator into a positive definite form, for which that paragraph gives the same limit. (P1) alone leaves \(t\in\{0,1,9,10\}\), and (P2) alone admits \(t=5\). The form \(h_0\) contracts the matrix indices and is not the emergent metric. The induced tensor of Definition~\ref{def:matrix-coords-ext} is its pull-back by a configuration, and for \(h_0\) of signature \((9,1)\) the restriction to a four-dimensional tangent space is positive definite, Lorentzian or degenerate according to the configuration, so the signature of the composite metric is a property of the background and is not derived here. (P1) is a condition on the representation of the integral within linear contours, and (P2) is supplied content, since Equation~\eqref{eq:Smatrix} is bosonic. This is no statement about non-linear contours, about linear maps that also mix the colour components, about whether the limit exists for \(2\le t\le8\), about fixed \(\gamma>0\), about \(N\to\infty\) taken first, or about Definitions~(B) and~(C). The signature selected under (P1) and (P2) is that of a contraction of matrix indices, both premises are imposed, and Assumption~\ref{ass:lorentzian} is not discharged.
